\documentclass[../../../include/open-logic-section]{subfiles}

\begin{document}

\olfileid{sfr}{infinite}{card-sb}

\olsection{પરિશિષ્ટ: શ્રેડર--બર્નસ્ટાઇનની સાબિતી}

સહજ ગણસિદ્ધાંત છોડીએ તે પહેલાં એક છેલ્લી સહજ (પણ સુવિકસિત!) સાબિતી
વિચારવાની બાકી છે. આ શ્રેડર--બર્નસ્ટાઇન
(\olref[sfr][siz][sb]{thm:schroder-bernstein})ની સાબિતી છે: જો
$\cardle{A}{B}$ અને $\cardle{B}{A}$ હોય, તો $\cardeq{A}{B}$; એટલે કે,
!!{injection}s $f \colon A \to B$ અને $g \colon B \to A$ આપેલાં હોય,
તો !!a{bijection} $h \colon A \to B$ હોય છે.

આ પ્રકરણમાં આપણે ડેડેકિન્ડના \emph{સંવરણો}ના ખ્યાલને અનુસર્યા છીએ.
ખરેખર, ડેડેકિન્ડે આ ખ્યાલ વાપરીને શ્રેડર--બર્નસ્ટાઇનની એક સુંદર
સાબિતી આપી હતી અને આપણે તેને અહીં રજૂ કરીશું. અહીંની સાબિતી
\citet[pp.~157--8]{Potter2004}ને નજીકથી અનુસરે છે; ત્યાં થોડું જુદું
પણ મૂળભૂત રીતે સરખું નિરૂપણ મળશે. થોડું ગૂગલમાં
શોધવાથી પણ ખાતરી થશે કે---$\sqrt{2}$ની અસંમેયતા જેવું---આ એવું પ્રમેય
છે જેની \emph{ઘણી} રસપ્રદ અને જુદી સાબિતીઓ છે.

\olref[sfr][infinite][dedekind]{Closure} જેવી જ સંજ્ઞા વાપરીને,
દરેક ગણ $B$ અને વિધેય~$f$ માટે લો:
\[
\Closureofunder{f}{B} = \bigcap \Setabs{X}{B \subseteq X
\text{ અને $X$ એ $f$-સંવૃત છે}}
\]
આ રીતે વ્યાખ્યાયિત $\Closureofunder{f}{B}$ એ $B$ને સમાવતો નાનામાં
નાનો $f$-સંવૃત ગણ છે, એ અર્થમાં કે:

\begin{lem}\ollabel{Closureprops}
	દરેક વિધેય $f$ અને દરેક $B$ માટે:
	\begin{enumerate}
		\item\ollabel{Closurehaselem} $B \subseteq \Closureofunder{f}{B}$; અને
		\item\ollabel{Closureclosed} $\Closureofunder{f}{B}$ એ $f$-સંવૃત છે; અને
		\item\ollabel{Closuresmallest} જો $X$ એ $f$-સંવૃત હોય અને $B
		\subseteq X$, તો $\Closureofunder{f}{B} \subseteq X$.
	\end{enumerate}
\end{lem}

\begin{proof}
બરાબર \olref[sfr][infinite][dedekind]{closureproperties}ની જેમ.
\end{proof}

શ્રેડર--બર્નસ્ટાઇન સુધી પહોંચવા માટે એક છેલ્લી હકીકત જોઈએ:

\begin{prop}\ollabel{sbhelper}
જો $A \subseteq B \subseteq C$ અને $A \approx C$ હોય, તો
$\cardeq{A}{B}$ અને $\cardeq{B}{C}$.
\end{prop}
\sourcecorrection{OLINF-002}{મૂળની \texttt{card-sb.tex}, પંક્તિઓ
50--52માં નિષ્કર્ષ $\cardeq{\cardeq{A}{B}}{C}$ તરીકે લખાયો છે,
જેમાં એક સભ્યસંખ્યા-સમાનતા બીજીના ગણ-સ્થાને ગેરરીતે આવેલી છે.
સાબિતી પહેલાં $C$થી $B$માં એક-એક વ્યાપ્ત વિધેય અને પછી $A$થી
$B$માં એવું વિધેય રચે છે; તેથી અહીં તે સાબિત કરતી બંને સુઘડ
સમાનતાઓ $\cardeq{A}{B}$ અને $\cardeq{B}{C}$ લખી છે.}

\begin{proof}
!!a{bijection} $f \colon C \to A$ આપેલું હોય, ત્યારે $F =
\Closureofunder{f}{C \setminus B}$ લો અને પ્રદેશ $C$ ધરાવતું વિધેય
$g$ આ પ્રમાણે વ્યાખ્યાયિત કરો:
\[
	g(x) =
	\begin{cases}
			f(x) &\text{જો $x \in F$}\\
			x & \text{અન્યથા}
		\end{cases}
\]
$g$ એ $C \to B$નું !!a{bijection} છે તે આપણે બતાવીશું; તેના પરથી
$\comp{f^{-1}}{g} \colon A \to B$ એ !!a{bijection} છે એવું મળશે અને
સાબિતી પૂરી થશે.

પહેલો દાવો એ છે કે જો $x \in F$ પણ $y\notin F$ હોય, તો
$g(x) \neq g(y)$. વિરોધાભાસ મેળવવા ઊલટું ધારો; ત્યારે
$y = g(y) = g(x) = f(x)$. \olref{Closureprops} મુજબ $F$ એ
$f$-સંવૃત છે અને $x \in F$ હોવાથી $y = f(x) \in F$, જે વિરોધાભાસ છે.

હવે ધારો કે $g(x) = g(y)$. તેથી ઉપરના દાવાથી $x \in F$ જો અને માત્ર
જો $y \in F$. જો $x, y \in F$, તો $f(x) = g (x) = g(y) = f(y)$;
અને $f$ એ !!a{bijection} હોવાથી $x = y$. જો $x, y \notin F$, તો
$x = g(x) = g(y) = y$. તેથી $g$ એ !!a{injection} છે.

હવે માત્ર $\ran{g} = B$ બતાવવાનું બાકી છે. તેથી $x \in B \subseteq C$
નિશ્ચિત કરો. જો $x \notin F$, તો $g(x) = x$. જો $x \in F$, તો કોઈ
$y \in F$ માટે $x = f(y)$; કારણ કે નહિતર $F \setminus \{x\}$ એ
$f$-સંવૃત હોત અને $C\setminus B$ને સમાવતો હોત, જે
\olref{Closureprops}થી અશક્ય છે. હવે $g(y) = f(y) = x$.
\end{proof}

છેલ્લે મુખ્ય પરિણામની સાબિતી આ રહી. યાદ કરો કે વિધેય $h$ અને ગણ $D$
આપેલાં હોય ત્યારે આપણે $\funimage{h}{D} = \Setabs{h(x)}{x
\in D}$ વ્યાખ્યાયિત કરીએ છીએ.

\begin{proof}[શ્રેડર--બર્નસ્ટાઇનની સાબિતી] ધારો કે $f \colon A \to B$
અને $g \colon B \to A$ એ !!{injection}s છે. $\funimage{f}{A}
\subseteq B$ હોવાથી $\funimage{g}{\funimage{f}{A}} \subseteq
g[B] \subseteq A$. વળી, $g$ અને $f$ બંને એક-એક હોવાથી
$\comp{f}{g} \colon A \to \funimage{g}{\funimage{f}{A}}$ પણ
!!{injection} છે; અને ખરેખર, તેનો સહપ્રદેશ જે રીતે વ્યાખ્યાયિત કર્યો
છે તે પરથી $\comp{f}{g}$ એ !!a{bijection} છે. તેથી
$\cardeq{\funimage{g}{\funimage{f}{A}}}{A}$ અને આમ
\olref{sbhelper} મુજબ કોઈ !!a{bijection} $h \colon A \to
\funimage{g}{B}$ છે. વળી, $g^{-1}$ એ !!a{bijection}
$\funimage{g}{B} \to B$ છે. તેથી $\comp{h}{g^{-1}} \colon A \to B$
એ !!a{bijection} છે.
\end{proof}

\end{document}
